\clearpage
\begin{summarybox}
\textbf{C089｜快记·快懂·快用·自检（师生共用）}

\textbf{一、快记：一个前提，三个结论。}
\[
\boxed{\begin{gathered}
P\in H,\ t_P\text{ 是 }P\text{ 处切线},\quad
A=t_P\cap\ell_1,\ B=t_P\cap\ell_2\\[2pt]
S_{\triangle AOB}=ab,\quad P=\frac{A+B}{2},\quad OA\cdot OB=a^2+b^2
\end{gathered}}
\]
还能马上得到 $S_{\triangle OAP}=S_{\triangle OPB}=ab/2$。
\textbf{逻辑方向只有}\quad$\text{切线}\Rightarrow\text{面积}=ab$，\textbf{不能倒推}。

\medskip
\textbf{二、快懂：为什么``一长一短''？}
\[
u=\frac{x_0}{a},\quad v=\frac{y_0}{b},\quad
u^2-v^2=1
\ \Rightarrow\ 
\underbrace{(u+v)}_{r}\underbrace{(u-v)}_{1/r}=1.
\]
\[
A=(ar,br),\qquad B=(a/r,-b/r).
\]
因此 $OA=|r|c,\ OB=c/|r|$（$c^2=a^2+b^2$），
而两渐近线的夹角固定且其正弦是 $2ab/c^2$，
所以 $S=\frac12\cdot c^2\cdot2ab/c^2=ab$。
\textbf{最小模型：}顶点处 $x=a$，底 $2b$、高 $a$，面积 $ab$。

\medskip
\textbf{三、快用：三条路线，按题目条件选择。}
\begin{enumerate}
\item \textbf{只问面积：}识别``切线＋两渐近线＋中心 $O$''，直接 $\boxed{ab}$。
\item \textbf{给切点：}用 $xx_0/a^2-yy_0/b^2=1$，
令 $u=x_0/a,v=y_0/b$；
$A=(a/(u-v),b/(u-v))$，$B=(a/(u+v),-b/(u+v))$。
\item \textbf{给直线 $y=kx+d$：}先验
$\boxed{d^2=a^2k^2-b^2>0}$，
再用 $x_Ax_B=a^2$ 求面积；竖直切线另作 $x=\pm a$。
\end{enumerate}
\textbf{坐标面积：}唯一正确的统一写法是
$\boxed{S_{\triangle OAB}=\frac12|x_Ay_B-x_By_A|}$。

\medskip
\textbf{四、学生自检（用来检验是否真正掌握）。}
\begin{enumerate}
\item $H:x^2/16-y^2/9=1$ 任意切线围出的面积？
\textbf{答：}$ab=4\cdot3=12$。
\item 上题若 $P$ 为切点，$S_{\triangle OAP}$ 多少？
\textbf{答：}$ab/2=6$，因为 $P$ 是 $AB$ 中点。
\item 若 \emph{任意直线} 与两渐近线构成面积 $12$，能推出它与 $H$ 相切吗？
\textbf{答：不能。}反例 $y=3$ 为割线，仍构成面积 $12$。
\end{enumerate}

\textbf{五、教师使用建议。}先用顶点图引出猜想；
让学生比较一组倍乘和倒数倍的具体坐标；
随后证明``切点是中点''，推出 $rs=1$；
最后介绍考场速证与斜率判别式。切忌将``中点性质''当作不证自明的图像结论。
\end{summarybox}
